\section{问题3：烘干时间确定}

\subsection{判据定义}

题面要求``药材\textbf{各处}的水分浓度应低于 $0.15$ kg/kg''，
``各处''即\textbf{全域}。故判据必须是
\begin{equation}
  \Cmax(t)=\max_{r\in[0,R_0]}C(r,t),
  \qquad
  \tstar=\inf\{\,t:\ \Cmax(t)<0.15\,\}.
\end{equation}

本文\textbf{不默认}最大值出现在中心，而是逐时刻对全域取 $\max$ 并记录
$\arg\max$ 的位置。实测确认 $\arg\max\equiv0$（中心），
即``中心是最难干燥的位置''是\textbf{被验证的结论，不是前提}。

\subsection{求解策略}

把 $\Cmax(t)-0.15$ 作为 ODE 求解器的\textbf{终止事件}，自适应精确定位零点。
两阶段：\textbf{粗扫}（按 60 s 输出、达标即停）$\to$
\textbf{二分}（在锚点与上界之间二分，每次试算都从\textbf{同一个固定锚点}
积分到试探时刻）。

\begin{table}[H]
  \centering
  \caption{问题3 主结果}
  \label{tab:p3main}
  \footnotesize
  \begin{tabular}{@{}ll@{}}
    \toprule
    项 & 值 \\
    \midrule
    阈值通过时刻 $\tstar$ & $\mathbf{206959.9521}$ s $=\mathbf{57.4889}$ h $=2.3954$ 天 \\
    交付取整（最早整数秒） & $206960$ s（\texttt{result3.xlsx} 末行） \\
    $\Cmax(\tstar)$ & $0.1499999998$（未舍入） \\
    最大值位置 & $r=0$（中心），逐时刻核对 \\
    二分次数 & 15 \\
    \bottomrule
  \end{tabular}
\end{table}

\textbf{注意}：$\Cmax(\tstar)=0.15$ 而非严格小于——这正是``连续模型的阈值交点''
的特征。本文报告的是\textbf{「阈值通过时刻」}，
并单独验证其后一段时间内全场持续达标。

\subsection{结果}

\begin{figure}[H]
  \centering
  \includegraphics[width=0.78\textwidth]{G6_p3_threshold.png}
  \caption{问题3 阈值通过时刻与终态含水率分布}
  \label{fig:G6}
\end{figure}

\begin{table}[H]
  \centering
  \caption{药材烘干过程的水分浓度（kg/kg）；$t_*$ = 57.4889 h}
  \label{tab:p3C}
  \footnotesize
  % 数据源：甲day2/results/M0/result3.xlsx；末行为达标时刻（不落在 60 s 网格上）
  \begin{tabular}{@{}lrrrrr@{}}
    \toprule
    时间/h & 0 cm & 0.5 cm & 1 cm & 1.5 cm & 2 cm \\
    \midrule
    6 & 1.0170 & 0.9882 & 0.9016 & 0.7551 & 0.5328 \\
    12 & 0.4562 & 0.4432 & 0.4031 & 0.3298 & 0.1642 \\
    18 & 0.2989 & 0.2916 & 0.2687 & 0.2253 & 0.0845 \\
    24 & 0.2378 & 0.2328 & 0.2168 & 0.1856 & 0.0665 \\
    30 & 0.2058 & 0.2019 & 0.1892 & 0.1641 & 0.0598 \\
    36 & 0.1859 & 0.1826 & 0.1718 & 0.1504 & 0.0566 \\
    42 & 0.1721 & 0.1692 & 0.1597 & 0.1407 & 0.0549 \\
    48 & 0.1618 & 0.1592 & 0.1507 & 0.1335 & 0.0537 \\
    54 & 0.1539 & 0.1515 & 0.1437 & 0.1277 & 0.0530 \\
    \textbf{烘干结束 57.4889 h} & 0.1500 & 0.1477 & 0.1402 & 0.1249 & 0.0527 \\
    \bottomrule
  \end{tabular}
\end{table}


完整结果见 \texttt{result3.xlsx}。

\begin{figure}[H]
  \centering
  \includegraphics[width=0.78\textwidth]{G12_max_center_surface.png}
  \caption{中心 / 表面 / 全域最大值 / 面积平均四条曲线——
           判据必须用 $\Cmax$ 而非平均值或表面值}
  \label{fig:G12}
\end{figure}

\subsection{分析与讨论}

\subsubsection*{（1）最大值出现在中心，必须用 $\Cmax$ 判定}

由 图 \ref{fig:G12} 可见，$\Cmax(t)$ 与中心 $C(0,t)$ \textbf{几乎重合}，
说明中心是最难干燥的位置。题面要求``药材各处的水分浓度应低于 0.15''，
故只有 $\Cmax$ 是合法判据。不同判据给出的达标时间：

\begin{table}[H]
  \centering
  \caption{三种判据的达标时间对照}
  \label{tab:criterion}
  \footnotesize
  \begin{tabular}{@{}lccc@{}}
    \toprule
    判据 & 达标时间 / h & 相对正确值的偏差 & 判定 \\
    \midrule
    $\Cmax(t)$ & $\mathbf{57.49}$ & --- & \textbf{正确} \\
    面积平均 $\bar C(t)$ & 36.04 & 低估 $21.44$ h（$\mathbf{37.3\%}$） & 错误 \\
    表面 $C(R_0,t)$ & 12.55 & 低估 $44.94$ h（$78.2\%$） & 错误 \\
    \bottomrule
  \end{tabular}
\end{table}

\begin{quote}
  \textbf{误用平均值会把 2.4 天的工艺说成 1.5 天。}
\end{quote}

\subsubsection*{（2）持续达标核验}

$\tstar$ 之后以 720 s 为间隔取 6 个探针点，$\Cmax$ 逐点核验
（表 \ref{tab:sustain}）：全部 $<0.15$，\textbf{持续达标}。
需要核验的理由是：若环境可能回湿，``首次达标''与``持续达标''不同。
本工况恒温段 $\Cinf\equiv0.05<0.15$，驱动势始终为正，预期不会回湿——
但本文\textbf{仍逐点核验，不靠预期}。

\begin{table}[H]
  \centering
  \caption{$\tstar$ 之后的持续达标核验}
  \label{tab:sustain}
  \footnotesize
  \begin{tabular}{@{}lccccc@{}}
    \toprule
    $t$ & $\tstar$ & $\tstar+720$ & $\tstar+1440$ & $\tstar+2160$ & $\tstar+3600$ \\
    \midrule
    $\Cmax$ & 0.1500000 & 0.1497910 & 0.1495835 & 0.1493772 & 0.1489688 \\
    \bottomrule
  \end{tabular}
\end{table}

\subsubsection*{（3）三类误差分开报告}

\begin{table}[H]
  \centering
  \caption{$\tstar$ 的误差三分量}
  \label{tab:err3}
  \footnotesize
  \begin{tabular}{@{}lll@{}}
    \toprule
    误差类型 & 量值 & 说明 \\
    \midrule
    事件定位误差 & $\pm0.0005$ s & 二分容差的一半 \\
    \textbf{空间离散误差} & $\approx\mathbf{17.4}$ \textbf{s} & \textbf{主导项}，由 $N$ 细化给出 \\
    时间积分误差 & 1.86 s & 一阶（Backward Euler），
                               拟合 $+1.29\sim+1.43$，\textbf{不声称二阶} \\
    \bottomrule
  \end{tabular}
\end{table}

\begin{quote}
  \textbf{定位到 1 s 不代表干燥时间整体准确到 1 s。}
  事件定位误差 0.5 ms 只是``在给定离散下找零点''的精度；
  真正的误差由空间离散与时间积分主导。
\end{quote}

\subsection{精度口径声明}
\label{sec:tstar-precision}

问题3 的答案依赖三个\textbf{必须显式声明}的口径选择。
换一个同样合理的取法，答案会移动几十到上千秒——\textbf{不声明，答案就不可比}。

\textbf{口径一：达标判据用全域最大值}（见 \S8.4.1）。
题面``各处''决定必须是 $\Cmax$；用面积平均会低估 $37.3\%$。

\textbf{口径二：恒温段环境取哪个窗口的均值}。
附件1 只覆盖 0--14400 s，$t>14400$ s 的边界必须假定：

\begin{table}[H]
  \centering
  \caption{环境平台窗口的取法对问题3 答案的影响}
  \label{tab:window}
  \footnotesize
  \begin{tabular}{@{}lccc@{}}
    \toprule
    平台窗口 & $T_\infty$ / $^\circ$C & $\Cinf$ / (kg/kg) & $\tstar$ / h \\
    \midrule
    $t\ge9600$ s（81 点，\textbf{本文采用}） & 50.0017 & 0.049998 & $\mathbf{57.4889}$ \\
    $t\ge12400$ s & 50.0092 & 0.049980 & 57.4747 \\
    直接取整数 & 50.0 & 0.05 & 57.4919 \\
    只取末点 & 50.1650 & 0.049860 & 57.1890 \\
    \bottomrule
  \end{tabular}
\end{table}

\textbf{跨度 $0.30$ h（$1090$ s）}——比本文的全部离散误差（约 $19$ s）
大 \textbf{57 倍}。故本文明确声明采用 $t\ge9600$ s 的实测均值。

\textbf{口径三：$\tstar$ 报告到几位小数}。
由表 \ref{tab:err3}，空间离散误差 $17.4$ s $\approx0.005$ h 落在第 3 位，
故本文报告 $\tstar\approx\mathbf{57.49}$ \textbf{h}（2 位小数）。
\textbf{题面要求四位小数的是表 5 的含水率数值，不是``烘干时间''这一时刻本身。}
把 $57.4889$ h 写进论文会让人误以为第 4 位有效——那是\textbf{虚假精度}。
